\section{Inductive bias as structural alignment}
\sectionmark{Inductive bias as structural alignment}

An inductive bias is the preference that makes a learning procedure choose some solutions over
other solutions that fit the available data
equally well.  It may enter through a prior, a regularizer, a function class, or the architecture
itself.  Battaglia et al. use this idea to argue that entity--relation structure should be built
into neural computation when the world is usefully described by entities, relations, and reusable
composition rules \cite{battaglia2018relational}.  A structural assumption that matches the
data-generating organization can
reduce sample complexity and support combinatorial generalization.

In concrete terms, choosing an inductive bias means writing a function whose organization can
correspond to regularities internal to the data.  That correspondence is a hypothesis about how
the problem is structured.  A representation $\phi_\theta$ chooses which
distinctions and relations become easy to express, while a function class $\mathcal H$ determines
which maps of those representations can be learned efficiently.  Let
\begin{equation}
m(x)=\mathbb E[Y\mid X=x]
\end{equation}
be the population response.  A structurally aligned representation--function pair makes
\begin{equation}
\inf_{f\in\mathcal H}
\mathbb E\!\left[\bigl(m(X)-f(\phi_\theta(X))\bigr)^2\right]
\label{eq:representation-bias-alignment}
\end{equation}
small while keeping $\mathcal H$ simple enough to estimate from the available data.  The
representation exposes a regularity; the function class is organized so that this regularity can
be expressed without relearning the structure from scratch.  This is the precise content of
``matching the function to the data structure.''  It is an architectural claim to be tested, not
a property guaranteed by scale, pretraining, or a graph-shaped input alone.

One-body structure is one special inductive bias, not a synonym for inductive bias.  An additive
class $f(x)=c+\sum_i h_i(x_i)$ prefers explanations assembled from independent site effects; a
pairwise Potts class enlarges that preference to include $G_{ij}(x_i,x_j)$; low-rank constraints
and sparse support impose further biases inside the pairwise class.  Other architectures may
instead prefer locality, symmetry, relational composition, hierarchy, or low-rank structure.
Whether one-body, pairwise, local, equivariant, compositional, or low-rank structure is appropriate
is an empirical claim about the endpoint in the chosen coordinates.

For coordinate models, symmetry supplies a particularly explicit example.  A global rigid motion
$g=(R,t)$ acts on every atomic coordinate by
\begin{equation}
g\cdot r_i=Rr_i+t,
\qquad R\in SO(3),\quad t\in\mathbb R^3,
\qquad g\in SE(3).
\label{eq:se3-action-on-coordinates}
\end{equation}
A coordinate-valued map is $SE(3)$-equivariant when
$F(g\cdot r)=g\cdot F(r)$; a scalar readout is invariant when
$q(g\cdot r)=q(r)$.  Restricting the function class to obey these identities builds the physical
irrelevance of the laboratory frame into the model before any data are observed.  This is a group-
theoretic inductive bias.  Reflections belong to the larger Euclidean group $E(3)$, but are not in
$SE(3)$; because protein stereochemistry is chiral, a mirror reflection is not generally a
physically equivalent coordinate change.

Finally, low approximation error on one benchmark establishes task alignment.  Alignment with the
scientific mechanism requires separating alternatives such as phylogeny,
experimental batch, database composition, or another nuisance regularity.  Prospective tests,
clade-aware splits, uncertainty calibration, and endpoint-specific interventions are what
distinguish useful structural alignment from shortcut alignment.

\subsection{Pallatom: a case study in structural alignment}

Pallatom makes this idea unusually concrete \cite{qu2025pallatom}.  Its coordinate convention,
relational architecture, denoising objective, and sequence readout are designed together around
one hypothesis: sequence identity, backbone shape, and side-chain packing can be generated
coherently by modeling a suitably represented all-atom coordinate distribution.

The first bias is introduced before the neural network sees any data.  A residue with variable
chemical topology is embedded into fourteen coordinate slots.  Native heavy atoms occupy their
standard slots; absent atoms are represented by virtual points at the $C_\alpha$ coordinate:
\begin{equation}
(x,a)\xmapsto{\Psi_{14}}
r_0\in\mathbb R^{N\times14\times3}.
\label{eq:pallatom-atom14-map}
\end{equation}
This map converts transdimensional generation into ordinary Euclidean point-cloud diffusion.  It
asserts that a residue is the stable compositional unit, that fourteen slots suffice for the
canonical amino acids, and that the artificial coincidence of unused slots at $C_\alpha$ is a
learnable nuisance.  These assumptions make the problem computationally regular; their status is
representational, while molecular cardinality and atom existence belong to the physical ontology.

The second bias is relational and multiscale in the sense emphasized by Battaglia et al.
\cite{battaglia2018relational}.  Pallatom uses local atom-level attention for fine geometry and a
globally communicating residue track with pair features for fold-scale organization.  Residue
features are broadcast to atoms; atom features are aggregated back to residues; denoised
coordinates are recycled into pair features; and triangle updates encourage relational
consistency.  The architecture therefore assumes
\begin{equation}
\text{local atomic geometry}
\longleftrightarrow
\text{residue identity and state}
\longleftrightarrow
\text{global residue--residue organization}.
\label{eq:pallatom-relational-bias}
\end{equation}
This hierarchy is well matched to folded proteins and avoids global attention over every atom
pair.  Long-range atom--atom effects are therefore mediated through a residue-level bottleneck,
which defines the model's relational capacity at fold scale.

The slogan ``learn $P(\text{all-atom})$'' should be read with care.  Pallatom supplies the denoiser
with a reduced side-chain labeling, and several amino acids share geometrically degenerate atom14
patterns.  Coordinates alone therefore leave sequence ambiguous.  The trained
object combines coordinate score matching with supervised sequence cross-entropy, residue- and
atom-level distogram losses, smooth-LDDT loss, and intermediate supervision.  Schematically,
\begin{equation}
\mathcal L
=\lambda\mathcal L_{\mathrm{atom}}
+\alpha_0\mathcal L_{\mathrm{seq}}
+\alpha_1\mathcal L_{\mathrm{LDDT}}
+\alpha_2\mathcal L_{\mathrm{dist,res}}
+\alpha_3\mathcal L_{\mathrm{dist,atom}}
+\alpha_4\mathcal L_{\mathrm{intermediate}}.
\label{eq:pallatom-composite-loss}
\end{equation}
Thus the representation, architecture, and auxiliary targets together make the sequence--structure
relation learnable beyond the information in an unsupervised coordinate density alone.
Pallatom encourages the $SE(3)$ identity statistically through coordinate centering, random
rigid-motion augmentation, and a Kabsch-aligned coordinate loss.  The function class still admits
departures from exact equivariance.  This distinction
is substantive: symmetry learned from augmentation is an approximate empirical regularity,
whereas symmetry imposed by the architecture holds by construction.

The training distribution contributes a third, unusually strong bias.  The reported model uses
short monomers, crops of length at most 128, 7,459 filtered PDB structures, and 27,697 heavily
curated AFDB structures.  AFDB candidates are required to have high predicted confidence and are
then filtered for packing density, core-residue fraction, loop content and length, radius of
gyration, secondary-structure segmentation, and redundancy.  Consequently the learned density is
better written as
\begin{equation}
p_\theta\!\left(r\mid r\in\mathcal F_{\mathrm{compact,designable}}
(\mathrm{PDB}\cup\mathrm{AFDB})\right)
\label{eq:pallatom-curated-distribution}
\end{equation}
than as a universal distribution over proteins.  This is a legitimate design choice: it aligns
the training set with the desired endpoint of compact, prediction-compatible monomers.  It also
limits claims about disordered proteins, complexes, alternate states, membrane environments, and
functional negative examples.

The ablations support this alignment interpretation.  In the reported setup, atom14 achieves
$87\%$ all-atom designability, compared with $5\%$ for a sequence-guided hybrid14 representation;
removing differentiable recycling lowers it to $21\%$.  These results identify the computational
representation that works with this denoiser and objective.  A claim about physical ontology would
require independent physical evidence.  Indeed, changing the sampling step scale raises reported
designability while reducing diversity and shifting the fraction of all-helical outputs.  The
sampler is therefore part of the inductive bias and changes which region of the learned
distribution is visited.

The evaluation closes the same loop.  Pallatom reports $85.03\%$ all-atom designability in its
primary comparison, where a sample passes when ESMFold assigns mean pLDDT above 80 and the
all-atom RMSD between the generated structure and the structure predicted from its generated
sequence is below $2$~\AA.  ProteinMPNN, ESMFold, Foldseek, and MMseqs2 provide useful tests of
self-consistency, backbone redesignability, novelty, and diversity.  Their own training
distributions and structural priors also enter the measurement.  High designability therefore
licenses the statement that the generated sequence and coordinates are mutually recoverable under
current predictors.  Stability, catalytic activity, binding specificity, equilibrium weights, and
experimental success remain separate validation endpoints.

Pallatom remains context-dependent even though its final sequence head is residuewise.  The head
pools atomic embeddings and applies a linear classifier after those embeddings have received global
residue, pair, triangle, and recycled-coordinate information.  The output logit at residue $i$ is
therefore a function of the generated molecular context.  A
local readout can therefore express pairwise and higher-order context already mixed into its input
representation.

Pallatom's main methodological achievement is therefore an alignment result.  A fixed-cardinality
all-atom representation, a local--global relational network, curated compact-protein data, a
composite geometric objective, and predictor-based designability metrics all point toward the
same target.  That alignment explains the model's strong in-silico performance more precisely
than saying that it simply ``learned chemistry.''  The remaining scientific question is whether
the aligned target is broad enough for the intended use; answering it requires prospective
synthesis and measurements under functional and negative-control conditions.

Structural alignment and conditional identifiability answer successive questions.  Pallatom shows
how a representation, architecture, objective, and evaluation loop can point coherently toward one
endpoint.  The next sections ask a different question: when an edge or router is interpreted as a
conditional relation, can its context be held fixed while the endpoint categories are varied?
That requirement leads to endpoint-removable context and exposes an information tradeoff that good
architectural alignment alone leaves unresolved.
